%! Permutations and Transposes — Unit 2, Lesson 2.4 — Warm-Up (Answer Key)
\documentclass[10pt]{article}
\usepackage{linalg-article}
\usepackage{linalg-key}   % pulls in -boxes; do NOT also load -boxes
\newcommand{\TallMath}[1]{$\displaystyle #1\rule[-1.4em]{0pt}{3.2em}$}
\newcommand{\xx}{\mathbf{x}}
\newcommand{\bb}{\mathbf{b}}
\newcommand{\tp}{\mathsf{T}}

\begin{document}

\pageheader{Unit 2, Lesson 2.4}{Warm-Up --- Answer Key}

\vspace{0.12in}

\begin{notesbox}{Getting Ready: Skills We'll Use Today}
\small Today we handle the one case Lesson 2.3 skipped --- a $0$ in the pivot --- by \emph{swapping rows},
and we meet the \emph{transpose}. These three quick items are the pieces.

\vspace{4pt}
\textbf{1. A swap matrix} (Lesson 1.3). Multiply, one entry at a time:
\[
  \begin{bmatrix} 0 & 1 \\ 1 & 0 \end{bmatrix}
  \begin{bmatrix} 4 \\ 9 \end{bmatrix}
  = \begin{bmatrix}{\color{keyred}\mathbf{9}}\\[2pt]{\color{keyred}\mathbf{4}}\end{bmatrix}.
\]
Notice this matrix simply \emph{swaps} the two entries.
\begin{work}
  \text{row 1} &= \text{second entry} = 9 \\
  \text{row 2} &= \text{first entry} = 4
\end{work}

\vspace{4pt}
\textbf{2. A stuck pivot} (Lesson 2.1/2.3). To start elimination on
$A=\begin{bsmallmatrix} 0&3\\2&5 \end{bsmallmatrix}$, the first pivot is the top-left entry. That entry is
${\color{keyred}\mathbf{0}}$, and elimination cannot start because we would have to divide by
{\color{keyred}\textbf{zero}}.

\vspace{4pt}
\textbf{3. Reading entries by address} (Lesson 1.3). For $A=\begin{bsmallmatrix} 6&1\\4&9 \end{bsmallmatrix}$,
the entry in \emph{row 1, column 2} is ${\color{keyred}\mathbf{1}}$ and the entry in \emph{row 2, column 1}
is ${\color{keyred}\mathbf{4}}$. (Today the transpose will trade these two spots.)
\end{notesbox}

\end{document}
